% Part: methods
% Chapter: induction
% Section: introduction

\documentclass[../../../include/open-logic-section]{subfiles}

\begin{document}

\olfileid{mth}{ind}{int}

\olsection{پېژندګلو}

استقرا د ثبوت يو مهم تخنيک دے چې په بېلابېلو بڼو کښې د منطق،
نظري کمپيوټرپوهنې او رياضياتو په تقريباً ټولو څانګو کښې کارېږي۔
د منطق د ډېرو نتيجو د ثابتولو دپاره ورته اړتيا وي۔

استقرا ډېر ځله د استنتاج په مقابله کښې يادېږي، او له ځانګړي څخه
عام ته د استدلال په توګه بيانېږي۔ د بېلګې په توګه، که ډېر شين
زمرد ووينو، او داسې څه و نۀ وينو چې زمرد ئې بولو خو شين نۀ وي،
نو ښايي دې نتيجې ته ورسېږو چې ټول زمرد شين دي۔ دا استقرايي
استدلال دے، ځکه چې له ډېرو ځانګړو حالتونو څخه، لکه دا زمرد شين
دے، هغه زمرد شين دے، او داسې نور، يوې عامې دعوې ته ځي چې ټول
زمرد شين دي۔ \emph{رياضي} استقرا هم داسې استدلال دے چې پايله ئې
يوه عامه دعوه وي، خو ډول ئې له دې «ساده استقرا» څخه ډېر بېل دے۔

په ډېر عمومي ډول، په رياضياتو کښې استقرايي ثبوت دې نتيجې ته
رسېږي چې د يو ټاکلي ډول ټول رياضيکي څيزونه يو ټاکلے خاصيت لري۔
په تر ټولو ساده حالت کښې هغه رياضيکي څيزونه چې استقرايي ثبوت
ورسره کار لري، طبيعي عددونه دي۔ په دې حالت کښې استقرايي ثبوت
دا ثابتوي چې ټول طبيعي عددونه يو خاصيت لري، او د دې دپاره ښيي چې
\begin{enumerate}
    \item $0$ هغه خاصيت لري، او
    \item هر کله چې يو عدد~$k$ هغه خاصيت ولري، نو~$k+1$ ئې هم لري۔
\end{enumerate}
پر طبيعي عددونو استقرا بيا ډېر ځله د داسې رياضيکي څيزونو په
هکله د عامو دعوو د ثابتولو دپاره هم کارېږي چې عددونه ورته
ټاکل کېداے شي۔ د بېلګې په توګه، هر متناهي سټ د غړو يو متناهي
شمېر~$n$ لري۔ که په استقرا وښيو چې هر عدد~$n$ دا خاصيت لري چې
«د~$n$ اندازې ټول متناهي سټونه \dots دي»، نو د ټولو متناهي
سټونو په هکله به مو يوه خبره ثابته کړې وي۔

استقرا هغو رياضيکي څيزونو ته هم تعميمېداے شي چې
\emph{په استقرايي ډول تعريف شوي} وي۔ د بېلګې په توګه، د يوې
رسمي ژبې عبارتونه، لکه د لومړي ترتيب د منطق عبارتونه، په
استقرايي ډول تعريفېږي۔ \emph{جوړښتي استقرا} د ټولو داسې
عبارتونو په هکله د نتيجو د ثابتولو يوه طريقه ده۔ په تېره بيا
جوړښتي استقرا په منطق کښې ډېره ګټوره ده او ډېر کارېږي۔

\end{document}
